\chapter{Parseo del DSL}
\label{ch:dsl}

El DSL es una pequeña gramática de cadena de flechas diseñada para
leerse en voz alta (``Un User se autentica a un Host y ejecuta un
Process'') y para entrar en una línea de input de UI. El parser
es un descendiente recursivo escrito a mano con una sola pasada
--- sin archivo de gramática, sin dependencia de generador.

\section{Gramática}

\begin{verbatim}
hypothesis      := step ( step )* modifier*
step            := entity_type arrow
arrow           := "-[" rel_type predicates? "]->" entity_type binding?
                   predicates? dest_predicates?
entity_type     := ident ( "(" binding ")" )? | "*"
rel_type        := ident | "*"
predicates      := "{" predicate ( "," predicate )* "}"
predicate       := key ("=" | "~" | " in ") (literal | list)
modifier        := having | within | k_clause
having          := "HAVING" agg_op "(" agg_arg ")" cmp literal
within          := "WITHIN" int time_unit
k_clause        := "{k=" int "}"
\end{verbatim}

Cada tipo de token se reconoce por una prueba de clase de
caracter, de modo que el parser hace exactamente una pasada
sobre el string de input. No hay backtracking; la gramática es
LL(1) tras el lookahead a nivel lexer para predicados.

\section{Algoritmo}

\begin{algorithm}
\caption{\textsc{Parse-Dsl}}\label{alg:parse-dsl}
\KwIn{string $\texttt{src}$, nombre opcional de hipótesis $\texttt{nm}$}
\KwOut{\code{Hypothesis} o \code{DslError}}
$\text{toks} \gets $ \textsc{Tokenize}($\texttt{src}$) \Comment*[r]{\Ord{|\texttt{src}|}}
$\text{steps} \gets \emptyset$\;
$\sigma_0 \gets $ \textsc{Parse-Entity-Type}(toks)\;
\While{toks tiene ``$\to$'' restante}{
    $\rho, P^{\text{edge}} \gets $ \textsc{Parse-Edge-Spec}(toks)\;
    $\sigma_i \gets $ \textsc{Parse-Entity-Type}(toks)\;
    $P^{\text{dst}} \gets $ \textsc{Parse-Dest-Predicates}(toks)\;
    empujar \code{HypothesisStep}$(\sigma_{i-1}, \rho, \sigma_i, P^{\text{edge}}, P^{\text{dst}})$ a steps\;
}
$k \gets$ parsear \texttt{\{k=N\}} opcional, default 1\;
$A \gets$ \textsc{Parse-Having}(toks)\;
$w \gets$ \textsc{Parse-Within}(toks)\;
\Return \code{Hypothesis}(nm, steps, $k$, $A$, $w$)\;
\end{algorithm}

\begin{complexity}
$\Ord{|\texttt{src}|}$ en tiempo y espacio. Cada token se visita
una sola vez; las listas de predicados basadas en \code{HashMap}
son $\Ord{1}$ para poblar. El espacio es lineal en la longitud de
la cadena parseada más los literales de predicados.
\end{complexity}

\begin{theorem}[Determinismo del parser]
\textsc{Parse-Dsl} es determinista: para todo input string fijo,
produce el mismo \code{Hypothesis} o el mismo \code{DslError} en
cada invocación.
\end{theorem}

\begin{proof}
No hay orden de iteración de hash-map, no hay valores aleatorios,
no hay llamadas al reloj del sistema en ningún lugar del parser.
Todas las estructuras construidas durante el parseo (las
\code{edge\_predicates} y \code{dest\_predicates} finales en
cada paso) preservan el orden de inserción. \qed
\end{proof}

\begin{inthecode}
\filepath{graph\_hunter\_core/src/dsl.rs:47} --- entry point
\code{parse\_dsl}.\\
\filepath{graph\_hunter\_core/src/dsl.rs} --- implementación de
gramática completa (helpers privados).
\end{inthecode}
